% Осень26!B3; нумерация внутри столбца без пустых строк.
% Источник: лекция 2, §16 (строка 399)
% Источник: лекция 2, §17 (строка 433)
% Источник: лекция 2, §19 (строка 488)
% Источник: лекция 2, §20 (строка 518)
\definitionticket{01}{Норма линейного оператора}
\label{definition:01}

\subsection*{Пространство линейных отображений. Норма оператора ---}

\textbf{Пространство \(\mathrm{Lin}(X,Y)\).} Для линейных пространств \(X,Y\) множество всех линейных отображений \(X\to Y\) обозначается \(\mathrm{Lin}(X,Y)\); в частности, \(\mathrm{Lin}(\mathbb{R}^m,\mathbb{R}^\ell)\). Это \textbf{линейное пространство}: элементы («операторы») реализуются матрицами, а матрица --- «просто длинющий вектор, записанный в виде таблички» \(\ell\times m\).

\begin{definition}

\emph{норма линейного оператора.} Для \(L\in\mathrm{Lin}(\mathbb{R}^m,\mathbb{R}^\ell)\) положим \[\|L\|:=\sup_{x\in\mathbb{R}^m,\ |x|=1}|Lx|,\] где \(|\cdot|\) --- евклидова норма в \(\mathbb{R}^m\) и \(\mathbb{R}^\ell\) (одна палочка), а \(\|\cdot\|\) --- \textbf{операторная} норма (две палочки) --- новый объект, «норма в пространстве линейных отображений». Пока это «временное обозначение» --- «через некоторое время повысим в статусе», проверив аксиомы нормы (§19).

Перебираем векторы единичной сферы, вычисляем \(Lx\), берём самый большой модуль.

\end{definition}

\textbf{Замечание 1 (конечность, эффективная оценка).} \(\|L\|<+\infty\): по \textbf{лемме об оценке нормы линейного отображения} (из 2 семестра): если \(L=(l_{ij})\) --- матрица оператора, то \[|Lx|\le C|x|,\qquad C=\sqrt{\sum_{i,j}l_{ij}^{2}},\] поэтому (при \(|x|=1\) все выражения под супремумом ограничены \(C\)) \[\|L\|\le\sqrt{\sum_{i,j}l_{ij}^{2}}.\]

\textbf{Замечание 2 (sup --- это max; теорема Вейерштрасса).} Единичная сфера \(S=\{x:|x|=1\}\) компактна в \(\mathbb{R}^m\) (замкнута и ограничена); отображение \(x\mapsto|Lx|\) непрерывно (композиция непрерывных). По теореме Вейерштрасса непрерывная функция на компакте ограничена и \textbf{достигает} наибольшего значения: \[\|L\|=\max_{|x|=1}|Lx|.\]

\textbf{Требует сверки.} \textbf{Контрольная точка (бесконечномерие).} Замена sup на max опирается на компактность сферы и \textbf{не переносится} на бесконечномерные нормированные пространства. Пример лектора: \(C[0,1]\) --- непрерывные функции на \([0,1]\) с нормой \(\|f\|=\max_{[0,1]}|f|\) (конечность --- та же теорема Вейерштрасса) --- бесконечномерное нормированное пространство; «точками» являются функции. В нём сфера \textbf{не} компактна: критерий «компактность \(\Leftrightarrow\) замкнутость + ограниченность» верен лишь в \(\mathbb{R}^m\); в бесконечномерном случае «это обычно рушится». Поэтому там sup на max заменять нельзя.

\textbf{Замечание 3 (главное неравенство).} Для всех \(x\in\mathbb{R}^m\): \[|Lx|\le\|L\|\cdot|x|.\]

\emph{Доказательство:} при \(x=0\) тривиально; при \(x\ne0\) вектор \(x/|x|\) единичный, поэтому \(|L(x/|x|)|\le\|L\|\), и по линейности \(|Lx|=|x|\cdot|L(x/|x|)|\le\|L\||x|\) («коэффициент растяжения выносится за знак нормы»).

\textbf{Замечание 4 (оптимальность константы).} Если найдётся \(C\) такое, что \[\forall x:\quad|Lx|\le C|x|,\] то \(\|L\|\le C\) (все выражения под супремумом по единичной сфере не превосходят \(C\)). Иными словами, \(\|L\|\) --- \textbf{наименьшая константа} в оценке \(|Lx|\le C|x|\): «самое точное неравенство получается при \(C=\|L\|\); меньше уже не возьмёте» (в \(\mathbb{R}^m\) максимум достигается --- существует \(x\), на котором равенство). В этом «востребованность» операторной нормы как технической величины для оценок.

\textbf{Геометрический смысл нормы} (пояснение лектора при доказательстве леммы §18): оператор переводит единичную сферу в некоторый «эллипсоид» --- сферу «раздуло, повернуло, продеформировало»; норма оператора --- \textbf{наибольший коэффициент растяжения}, «какое место этой сферы дальше всего отъехало от нуля» (оператор переводит \(0\) в \(0\)).

\subsection*{Примеры вычисления нормы ---}

\begin{enumerate}
\def\labelenumi{\arabic{enumi}.}
\tightlist
\item
  \textbf{\(m=\ell=1\):} \(L:\mathbb{R}\to\mathbb{R}\), \(Lx=l_{11}x\) (матрица \(1\times1\)). Единичная «сфера» --- два точки \(\{\pm1\}\), и \[\|L\|=|l_{11}|.\]
\item
  \textbf{\(m=1\), \(\ell\) любое:} \(L:\mathbb{R}\to\mathbb{R}^\ell\); общий вид линейного отображения --- \(Lx=x\cdot v\), где \(v\in\mathbb{R}^\ell\) --- вектор-столбец (матрица \(\ell\times1\) и «физически является вектором»). При \(|x|=1\) (т.е. \(x=\pm1\)): \(|Lx|=|v|\), поэтому \[\|L\|=|v|.\]
\item
  \textbf{\(\ell=1\), \(m\) любое:} \(L:\mathbb{R}^m\to\mathbb{R}\); общий вид --- \(Lx=\langle v,x\rangle\) (скалярное произведение; матрица --- строка \(v\)). По неравенству Коши--Буняковского \[|\langle v,x\rangle|\le|v|\cdot|x|=|v|\quad\text{при }|x|=1,\] и равенство достигается, когда \(x\) сонаправлен \(v\) (угол \(0\) или \(\pi\), косинус \(\pm1\)). Поэтому \[\|L\|=|v|.\]
\item
  \textbf{Общий случай \(m,\ell\):} явной формулы нет, норма «плохо считается». Комментарий лектора: евклидова норма матрицы \(\sqrt{\sum l_{ij}^2}\) «хорошо считается, но даёт завышенное представление о том, как действует оператор» и в оценках оказывается малоэффективной; операторная норма, наоборот, «плохо считается, зато все оценки с ней на ура самые лучшие». «Человечество долго думало, что лучше» --- выбор в пользу операторной нормы.
\end{enumerate}

\textbf{Обобщение (произвольные нормированные пространства).} Всё определение переносится на случай, когда \(X,Y\) --- \textbf{любые нормированные пространства} (нужны лишь: норма в \(X\) --- чтобы брать единичные векторы, и норма в \(Y\) --- чтобы измерять образы): \[\|A\|=\sup_{\|x\|_X=1}\|Ax\|_Y .\] В конечномерном случае это максимум; в бесконечномерном --- вообще говоря, только супремум (и его конечность не гарантирована).

\subsection*{\texorpdfstring{Теорема: \(L\mapsto\|L\|\) --- норма на \(\mathrm{Lin}(X,Y)\); композиция ---}{Теорема: L\textbackslash mapsto\textbackslash\textbar L\textbackslash\textbar{} --- норма на \textbackslash mathrm\{Lin\}(X,Y); композиция ---}}

\begin{theorem}

Отображение \(\mathrm{Lin}(X,Y)\to\mathbb{R}\), \(L\mapsto\|L\|\), --- \textbf{норма} в линейном пространстве операторов:

\textbf{(a)} \(\|L\|\ge0\), причём \(\|L\|=0\iff L\) --- нулевой оператор; \textbf{(b)} положительная однородность: \(\|\alpha L\|=|\alpha|\cdot\|L\|\) для всех \(\alpha\in\mathbb{R}\); \textbf{(c)} неравенство треугольника: \(\|L+M\|\le\|L\|+\|M\|\).

\end{theorem}

\textbf{Дополнительно (мультипликативность).} Если \(A\in\mathrm{Lin}(X,Y)\), \(B\in\mathrm{Lin}(Y,Z)\) (\(Z\) --- нормованное пространство), то композиция \(BA\in\mathrm{Lin}(X,Z)\) и \[\|BA\|\le\|B\|\cdot\|A\|.\]

«Теорема, которая наконец узаконит название ``норма \(L\)»''.

\begin{proof}

\textbf{(a)} Неотрицательность очевидна (супремум неотрицательных величин). Если \(\|L\|=0\), то \(Lx=0\) в \textbf{каждой} точке единичной сферы, а тогда по линейности \(L=0\) всюду (любой \(x\ne0\) --- растянутый единичный вектор). Обратно, у нулевого оператора норма \(0\).

\textbf{(b)} «Домножьте оператор на \(\alpha\) --- весь супремум домножится на \(|\alpha|\)»: \(\|\alpha L\|=\sup_{\|x\|=1}\|\alpha Lx\|=|\alpha|\sup_{\|x\|=1}\|Lx\|=|\alpha|\|L\|\) (обычное свойство супремума при растяжении множества значений; однородности нормы в \(Y\)).

\textbf{(c)} Для единичного \(x\) (\(\|x\|_X=1\)): \[\|(L+M)x\|_Y\le\|Lx\|_Y+\|Mx\|_Y\le\|L\|+\|M\|\] --- сначала неравенство треугольника для нормы в \(Y\) (обычные векторы!), затем замечание 3 из §16 (\(\|Lx\|\le\|L\|\), \(\|Mx\|\le\|M\|\)). Оценка верна при всех единичных \(x\), значит и для супремума: \[\|L+M\|\le\|L\|+\|M\|.\]

\end{proof}

\textbf{Композиция.} Для всех \(x\): \[\|(BA)x\|_Z=\|B(Ax)\|_Z\le\|B\|\cdot\|Ax\|_Y\le\|B\|\cdot\|A\|\cdot\|x\|_X\] (дважды замечание 3). Значит, \(C=\|B\|\|A\|\) --- константа в оценке из замечания 4, откуда \(\|BA\|\le\|B\|\|A\|\). (Эквивалентно: sup по единичной сфере.) \(\blacksquare\)

\subsection*{Сводные формулы для нормы оператора ---}

Итоговое замечание лектора: в \(\mathrm{Lin}(X,Y)\) норму оператора \(A\) можно вычислять по любой из формул: \[
\begin{aligned}
\|A\|&=\sup_{\|x\|_X=1}\|Ax\|_Y
      =\sup_{\|x\|_X\le1}\|Ax\|_Y\\
 &=\sup_{x\ne0}\frac{\|Ax\|_Y}{\|x\|_X}\\
 &=\inf\{C:\ \forall x\quad\|Ax\|_Y\le C\|x\|_X\}.
\end{aligned}
\]

Пояснения лектора:

\begin{itemize}
\tightlist
\item
  по \textbf{открытому} единичному шару sup тот же: «сколь угодно близко к точкам сферы подойдёте --- значение подойдёт к нужному» (доказательство --- упражнение); можно и по замкнутому шару;
\item
  по всему пространству --- с оговоркой \(x\ne0\) (нормировка \(x/|x|\) «вносит единицу нормы под числитель»);
\item
  последняя формула --- «суперпродвинутая»: перебрать все константы \(C\), для которых верна оценка из замечания 4, и взять наименьшую.
\end{itemize}

\textbf{Требует сверки.} \textbf{Про достижимость inf/max:} в конечномерном пространстве теорема Вейерштрасса «нас от всего спасёт» --- супремумы и инфимумы достигаются. В бесконечномерном это отдельный вопрос («есть такой вид спорта --- переписывать инфимумы в минимумы»); обычно «это никого не тревожит», но осознавать, что достижимость надо выяснять, --- полезно.


В формуле с единичным шаром можно брать как замкнутый, так и открытый шар.
